% GENERATED by paper/gen_prompt_appendix.py on 2026-09-30 17:05Z from prompts/noilai.yaml + prompts/demos.yaml; do not edit.
% prompt file sha256: demos.yaml=5fe710f99645, noilai.yaml=a390e6a44153, xcopa.yaml=de275fa2df4b, _all=002fb57ab3d6
\paragraph{Transformation (\task{T1}), variant \var{1}, paraphrase p0, three demonstrations.} Prompt id \texttt{vi-p0-s3-explained-raw}, sha256 \texttt{fc90f3c2182f}.
\begin{quote}\small
Nói lái là một cách chơi chữ trong tiếng Việt: hai âm tiết của một cụm từ trao đổi với nhau một số thành phần (phụ âm đầu, vần hoặc thanh điệu) để tạo thành một cụm từ mới. \\
Mỗi âm tiết tiếng Việt gồm ba phần: phụ âm đầu (có thể không có), vần (âm đệm nếu có, âm chính và âm cuối nếu có) và thanh điệu (một trong sáu thanh: ngang, huyền, sắc, hỏi, ngã, nặng). Ví dụ: \textquotedbl{}trung\textquotedbl{} gồm phụ âm đầu \textquotedbl{}tr\textquotedbl{}, vần \textquotedbl{}ung\textquotedbl{} và thanh ngang; \textquotedbl{}bình\textquotedbl{} gồm phụ âm đầu \textquotedbl{}b\textquotedbl{}, vần \textquotedbl{}inh\textquotedbl{} và thanh huyền.

Kiểu nói lái cần áp dụng: đổi vần, giữ phụ âm đầu và thanh điệu. \\
Cách làm từng bước: \\
1. Tách mỗi âm tiết thành phụ âm đầu, vần và thanh điệu. \\
2. Đổi vần của âm tiết thứ nhất cho vần của âm tiết thứ hai. \\
3. Giữ nguyên phụ âm đầu và thanh điệu của mỗi âm tiết ở vị trí cũ. \\
4. Ghép lại thành hai âm tiết và viết đúng chính tả. \\
Khi ghép lại, hãy viết đúng chính tả tiếng Việt: âm /k/ viết là \textquotedbl{}k\textquotedbl{} trước e, ê, i, viết là \textquotedbl{}q\textquotedbl{} trước âm đệm u và viết là \textquotedbl{}c\textquotedbl{} trong các trường hợp còn lại (ví dụ \textquotedbl{}c\textquotedbl{} ghép với vần \textquotedbl{}eo\textquotedbl{} phải viết là \textquotedbl{}keo\textquotedbl{}, không viết \textquotedbl{}ceo\textquotedbl{}); tương tự, \textquotedbl{}gh\textquotedbl{} và \textquotedbl{}ngh\textquotedbl{} đứng trước e, ê, i (ví dụ \textquotedbl{}nghe\textquotedbl{}, không viết \textquotedbl{}nge\textquotedbl{}); dấu thanh đặt đúng vị trí.

Ví dụ: \\
Cụm từ: trung bình \\
Đáp án: trinh bùng

Cụm từ: làm chủ \\
Đáp án: lù chảm

Cụm từ: vô hình \\
Đáp án: vinh hồ

Bây giờ hãy áp dụng kiểu nói lái trên cho cụm từ sau. \\
Cụm từ: mèo cái \\
Chỉ viết cụm từ kết quả trên một dòng riêng ở cuối, theo đúng dạng: \\
Đáp án: \textless{}cụm từ kết quả\textgreater{}
\end{quote}
\paragraph{Decoding (\task{T2}), paraphrase p0, three demonstrations (the variant is withheld).} Prompt id \texttt{vi-p0-s3-explained-raw}, sha256 \texttt{785246b4733d}.
\begin{quote}\small
Nói lái là một cách chơi chữ trong tiếng Việt: hai âm tiết của một cụm từ trao đổi với nhau một số thành phần (phụ âm đầu, vần hoặc thanh điệu) để tạo thành một cụm từ mới. Có bốn kiểu nói lái thường gặp: \\
- đổi vần, giữ phụ âm đầu và thanh điệu: Đổi vần của âm tiết thứ nhất cho vần của âm tiết thứ hai. Giữ nguyên phụ âm đầu và thanh điệu của mỗi âm tiết ở vị trí cũ. \\
- đổi cả phụ âm đầu và vần, giữ thanh điệu: Đổi chỗ phần phụ âm đầu và vần của hai âm tiết cho nhau (tức là đổi chỗ hai âm tiết). Giữ thanh điệu ở vị trí cũ: âm tiết đứng trước vẫn mang thanh của âm tiết đứng trước, âm tiết đứng sau vẫn mang thanh của âm tiết đứng sau. \\
- đổi thanh điệu, giữ phụ âm đầu và vần: Đổi thanh điệu của âm tiết thứ nhất cho thanh điệu của âm tiết thứ hai. Giữ nguyên phụ âm đầu và vần của mỗi âm tiết. \\
- đổi vần và thanh điệu, giữ phụ âm đầu: Đổi cả vần và thanh điệu của âm tiết thứ nhất cho vần và thanh điệu của âm tiết thứ hai. Giữ nguyên phụ âm đầu của mỗi âm tiết ở vị trí cũ. \\
Mỗi âm tiết tiếng Việt gồm ba phần: phụ âm đầu (có thể không có), vần (âm đệm nếu có, âm chính và âm cuối nếu có) và thanh điệu (một trong sáu thanh: ngang, huyền, sắc, hỏi, ngã, nặng). Ví dụ: \textquotedbl{}trung\textquotedbl{} gồm phụ âm đầu \textquotedbl{}tr\textquotedbl{}, vần \textquotedbl{}ung\textquotedbl{} và thanh ngang; \textquotedbl{}bình\textquotedbl{} gồm phụ âm đầu \textquotedbl{}b\textquotedbl{}, vần \textquotedbl{}inh\textquotedbl{} và thanh huyền.

Ví dụ: \\
Cụm từ: trinh bùng \\
Đáp án: trung bình

Cụm từ: chù lảm \\
Đáp án: làm chủ

Cụm từ: vồ hinh \\
Đáp án: vô hình

Cụm từ sau đây là một cách nói lái. Cụm từ gốc là gì? \\
Cụm từ: mài kéo \\
Chỉ viết cụm từ gốc trên một dòng riêng ở cuối, theo đúng dạng: \\
Đáp án: \textless{}cụm từ gốc\textgreater{}
\end{quote}
\paragraph{Validity (\task{T3}), variant \var{1}, paraphrase p0, a spelling twin as candidate.} Prompt id \texttt{vi-p0-s3-explained-raw}, sha256 \texttt{855d848b82af}.
\begin{quote}\small
Nói lái là một cách chơi chữ trong tiếng Việt: hai âm tiết của một cụm từ trao đổi với nhau một số thành phần (phụ âm đầu, vần hoặc thanh điệu) để tạo thành một cụm từ mới. \\
Mỗi âm tiết tiếng Việt gồm ba phần: phụ âm đầu (có thể không có), vần (âm đệm nếu có, âm chính và âm cuối nếu có) và thanh điệu (một trong sáu thanh: ngang, huyền, sắc, hỏi, ngã, nặng). Ví dụ: \textquotedbl{}trung\textquotedbl{} gồm phụ âm đầu \textquotedbl{}tr\textquotedbl{}, vần \textquotedbl{}ung\textquotedbl{} và thanh ngang; \textquotedbl{}bình\textquotedbl{} gồm phụ âm đầu \textquotedbl{}b\textquotedbl{}, vần \textquotedbl{}inh\textquotedbl{} và thanh huyền.

Kiểu nói lái đang xét: đổi vần, giữ phụ âm đầu và thanh điệu. \\
Cách làm từng bước: \\
1. Tách mỗi âm tiết thành phụ âm đầu, vần và thanh điệu. \\
2. Đổi vần của âm tiết thứ nhất cho vần của âm tiết thứ hai. \\
3. Giữ nguyên phụ âm đầu và thanh điệu của mỗi âm tiết ở vị trí cũ. \\
4. Ghép lại thành hai âm tiết và viết đúng chính tả. \\
Khi ghép lại, hãy viết đúng chính tả tiếng Việt: âm /k/ viết là \textquotedbl{}k\textquotedbl{} trước e, ê, i, viết là \textquotedbl{}q\textquotedbl{} trước âm đệm u và viết là \textquotedbl{}c\textquotedbl{} trong các trường hợp còn lại (ví dụ \textquotedbl{}c\textquotedbl{} ghép với vần \textquotedbl{}eo\textquotedbl{} phải viết là \textquotedbl{}keo\textquotedbl{}, không viết \textquotedbl{}ceo\textquotedbl{}); tương tự, \textquotedbl{}gh\textquotedbl{} và \textquotedbl{}ngh\textquotedbl{} đứng trước e, ê, i (ví dụ \textquotedbl{}nghe\textquotedbl{}, không viết \textquotedbl{}nge\textquotedbl{}); dấu thanh đặt đúng vị trí.

Ví dụ: \\
Cụm từ gốc: trung bình; cụm từ cần xét: trinh bùng \\
Đáp án: Có

Cụm từ gốc: làm chủ; cụm từ cần xét: chù lảm \\
Đáp án: Không

Cụm từ gốc: vô hình; cụm từ cần xét: vinh hồ \\
Đáp án: Có

Câu hỏi: Cụm từ \textquotedbl{}mài céo\textquotedbl{} có phải là kết quả đúng khi nói lái cụm từ \textquotedbl{}mèo cái\textquotedbl{} theo kiểu trên không? Kết quả đúng phải đúng cả phụ âm đầu, vần, thanh điệu và chính tả. \\
Trả lời \textquotedbl{}Có\textquotedbl{} hoặc \textquotedbl{}Không\textquotedbl{} trên một dòng riêng ở cuối, theo đúng dạng: \\
Đáp án: \textless{}Có hoặc Không\textgreater{}
\end{quote}
